\documentclass[10pt,a4paper]{amsart}
\usepackage[margin=19mm]{geometry}
\usepackage{amsmath,amssymb,mathtools}
\usepackage[scr=rsfs,cal=euler]{mathalfa}
\usepackage{xfrac}
\usepackage[unicode,hidelinks,bookmarks=false]{hyperref}
\newcommand{\ie}{i.e.\ }
\newcommand{\cf}{cf.\ }
\newcommand{\resp}{respectively\ }
\newcommand{\Subsection}{\subsection}
\providecommand{\nobreakdash}{\nobreak\hbox{-}}
\setlength{\emergencystretch}{2em}
\multlinegap0pt
\usepackage{bm}
\usepackage{bbm}
\usepackage{dsfont}
\usepackage{xspace}
\usepackage{cancel}
\usepackage{mathtools}
\usepackage{amsthm}
\makeatletter
\@ifundefined{thm}{\newtheorem{thm}{Thm}}{}
\@ifundefined{lem}{\newtheorem{lem}[thm]{Lemma}}{}
\@ifundefined{prop}{\newtheorem{prop}[thm]{Proposition}}{}
\makeatother
\newcommand{\HS}{\mathrm{HS}}
\newcommand{\II}{\mathrm{II}}
\newcommand{\cH}{\mathcal{H}}
\newcommand{\dd}{\,d}
\title[On the Brunn-Minkowski inequality for $q$-th dual quermassin]{On the Brunn-Minkowski inequality for $q$-th dual quermassintegrals with $q>n$}
\author{Haizhong Li, Yao Wan, Yijia Zhang}
\date{Source: arXiv:2608.03949v1 (CC BY 4.0)}
\begin{document}
\maketitle

\noindent\emph{Source and licence.} On the Brunn-Minkowski inequality for $q$-th dual quermassintegrals with $q>n$. Version arXiv:2608.03949v1 (CC BY 4.0), distributed under the Creative Commons Attribution 4.0 International licence, as recorded in the arXiv licence metadata for this exact version and shown on the abstract page https://arxiv.org/abs/2608.03949v1. Full rights notice: licenses/arxiv-2608.03949-CC-BY-4.0.txt.\par\smallskip

\noindent\textit{Editorial scope.} Counted unit: the Prop whose source label is \texttt{prop-local-logconcave-strict} in the article (source0.tex lines 1695--1702, source label \texttt{prop-local-logconcave-strict}), delivered together with the article's own complete original proof of it (source0.tex lines 1704--1775, one \texttt{proof} environment of 72 lines). No mathematics is written, repaired or re-derived: statement and proof are the article's own text line for line. Required context retained uncounted: lem-variation (lines 279--291); lem-neumann (lines 895--958); eq-reilly-final (lines 1308--1324); prop-energy (lines 1603--1609). Every definition, display and label of those lines is retained unchanged. Omitted from this package and not redistributed: 1--278, 292--894, 959--1307, 1325--1602, 1610--1694, 1703--1703, 1776--2335. Nothing inside a retained excerpt is omitted or altered: every retained body is delivered exactly as the article's lines are written, so no omission receipt of any kind accompanies this package. Citations: the retained excerpts carry no citation command at all, so no bibliography entry of the article is transcribed and no reference is redistributed. Reference adaptation: none was needed and none was made; every reference inside the delivered unit resolves inside it, so no unresolved reference or citation is produced. Printed slips kept literal: no printed word, symbol, hypothesis or number of the counted statement or of its proof was altered. Layout and class: the article's own class is \texttt{amsart} with packages including enumerate, amsrefs, amsfonts, amsmath, amssymb, amscd, amsthm, bm, cancel, url, graphicx, hyperref, multicol, comment; the wrapper is the lane's pinned \texttt{amsart} editorial wrapper at 10pt on A4, which restates the theorem-like environments the retained text needs and adds the article's own macro definitions that the retained text needs (\texttt{\textbackslash HS}, \texttt{\textbackslash II}, \texttt{\textbackslash cH}, \texttt{\textbackslash dd}), with extra wrapper packages (bm, bbm, dsfont, xspace, cancel, mathtools). The delivered numbering is the unit's own. Assets: no retained span contains a figure environment, a graphics-inclusion command, a PDF-page inclusion, a table or a TikZ picture; no image file is copied into this package, the package declares no asset dependency and no image file is redistributed with it. Licence: version arXiv:2608.03949v1 of this article is distributed under the Creative Commons Attribution 4.0 International licence, as recorded in the arXiv licence metadata for this exact version and shown on the abstract page https://arxiv.org/abs/2608.03949v1. A full rights notice is delivered at licenses/arxiv-2608.03949-CC-BY-4.0.txt. Characters: every retained excerpt is pure ASCII, so the delivered engine, XeLaTeX with its default Latin Modern text font, reports no missing character.
\begin{lem}[Weighted first and second variation]\label{lem-variation}
For the above perturbation,
\begin{align}
    \left.\frac{d}{dt}\right|_{t=0}\mu_q(K_t)
    &=\int_{\partial K}\phi\,d\mu_{\partial K,q},\\
    \left.\frac{d^2}{dt^2}\right|_{t=0}\mu_q(K_t)
    &=\int_{\partial K}\left(H_W\phi^2-\left\langle \II^{-1}\nabla_{\partial K}\phi,\nabla_{\partial K}\phi\right\rangle\right)d\mu_{\partial K,q},
\end{align}
where
\[
    d\mu_{\partial K,q}=|x|^{q-n}d\cH^{n-1}_{\partial K}.
\]
\end{lem}

\begin{lem}\label{lem-neumann}
  Let $-n<\alpha<n$, and let $K\in\mathcal{K}_{+}^n$. Let
  $\varphi\in C^\infty(\partial K)$ and set
  \[
    \dd\tilde{\mu}_\alpha=|x|^\alpha\,\dd x,
    \qquad
    \dd\tilde{\mu}_{\alpha,\partial K}
    =
    |x|^\alpha\,\dd H^{n-1},
  \]
  and
  \[
    C_\varphi
    =
    \frac{
      \int_{\partial K}\varphi\,\dd\tilde{\mu}_{\alpha,\partial K}}
    {      \tilde{\mu}_\alpha(K)}.
  \]
  Then the normalized weighted Neumann problem
  \[
    L_\alpha u=C_\varphi\quad\text{in }K,
    \qquad
    u_\nu=\varphi\quad\text{on }\partial K,
    \qquad
    \int_K u\,\dd\tilde{\mu}_\alpha=0
  \]
  has a unique weak solution
  $u\in H^1_\diamond(K,\tilde{\mu}_\alpha)$, where
  \[
    H^1_\diamond(K,\tilde{\mu}_\alpha)
    =
    \left\{
      v\in H^1(K,\tilde{\mu}_\alpha):
      \int_K v\,\dd\tilde{\mu}_\alpha=0
    \right\}.
  \]
  Equivalently, $u$ is the unique element of
  $H^1_\diamond(K,\tilde{\mu}_\alpha)$ satisfying
  \begin{equation}\label{eq-neumann-weak}
    \int_K\langle\nabla u,\nabla\eta\rangle\,
    \dd\tilde{\mu}_\alpha
    =
    -C_\varphi\int_K\eta\,\dd\tilde{\mu}_\alpha
    +
    \int_{\partial K}\varphi\eta\,
    \dd\tilde{\mu}_{\alpha,\partial K}
  \end{equation}
  for every $\eta\in H^1(K,\tilde{\mu}_\alpha)$.

    Moreover,
  \[
    u\in C^\infty(\operatorname{int}K\setminus\{0\}),
    \qquad
    L_\alpha u=C_\varphi\quad\text{in }K\setminus\{0\},
    \qquad
    u_\nu=\varphi\quad\text{on }\partial K,
  \]
  and, for every $0<\varepsilon<\operatorname{dist}(0,\partial K)$,
  \[
    u\in W^{2,2}\bigl(K\setminus\overline{B_\varepsilon}\bigr).
  \]
  If in addition $K$ and $\varphi$ are unconditional, then $u$ is
  unconditional.
\end{lem}

  \begin{equation}\label{eq-reilly-final}
    \begin{aligned}
      C_\varphi^2\tilde{\mu}_\alpha(K)
      =&\int_K\bigl(\|D^2u\|_{\HS}^2
         +D^2W_\alpha(\nabla u,\nabla u)\bigr)
         \dd\tilde{\mu}_\alpha
       +\int_{\partial K}H_{W_\alpha}u_\nu^2
         \,\dd\tilde{\mu}_{\alpha,\partial K}\\
       &+\int_{\partial K}\mathrm{II}(\nabla_{\partial K}u,
         \nabla_{\partial K}u)
         \,\dd\tilde{\mu}_{\alpha,\partial K}
       -2\int_{\partial K}
         \langle\nabla_{\partial K}u_\nu,
         \nabla_{\partial K}u\rangle
         \,\dd\tilde{\mu}_{\alpha,\partial K}.
    \end{aligned}
  \end{equation}

\begin{prop}\label{prop-energy}
  Let $K\in\mathcal{K}_{+}^n$ be unconditional, and let $u$ be
  the unconditional Neumann solution from Lemma~\ref{lem-neumann}. If
  $-n<\alpha\le1$, then $I_\alpha(u)\ge0$.
  Moreover, if $I_\alpha(u)=0$, then $u$ is constant on $K$, and
  consequently $\varphi = u_\nu = 0$ on $\partial K$.
\end{prop}

\begin{prop}\label{prop-local-logconcave-strict} 
  Under the above assumptions,
  \[
    \left.\frac{\dd^2}{\dd s^2}\right|_{s=0}
    \log \tilde{\mu}(s)\le 0 .
  \]
  Moreover, if $\psi\not\equiv 0$, the inequality is strict.
\end{prop}

\begin{proof}
  Since $\psi\in C^2_{und}(\mathbb{S}^{n-1})$ and $\nu_K$ is $C^1$, we
  have $\varphi\in C^1(\partial K)$. As noted after
  Lemma~\ref{lem-neumann}, the Neumann solution exists and satisfies
  the same conclusions as in the smooth case.

  It suffices to show $\tilde{\mu}(0)\tilde{\mu}''(0)\le\tilde{\mu}'(0)^2$.
  By Lemma~\ref{lem-variation},
  \[
    \tilde{\mu}(0)=\tilde{\mu}_\alpha(K),\quad\tilde{\mu}'(0)=C_\varphi\tilde{\mu}_\alpha(K),
  \]
  and
  \[
  \tilde{\mu}''(0)=\int_{\partial K}\bigl(H_{W_\alpha}\varphi^2-\mathrm{II}^{-1}(\nabla_{\partial K}\varphi,\nabla_{\partial K}\varphi)\bigr)\,d\tilde{\mu}_{\alpha,\partial K}.
  \]
  Applying the singular weighted Reilly formula \eqref{eq-reilly-final} and Proposition \ref{prop-energy}, we obtain
   \begin{align*}
    \frac{\tilde{\mu}'(0)^2}{\tilde{\mu}(0)}
    &=
    I_\alpha(u)
    +
    \int_{\partial K}
      H_{W_\alpha}\varphi^2\,
      \dd\tilde{\mu}_{\alpha,\partial K}  \\
    &\quad
    +
    \int_{\partial K}
      \mathrm{II}(\nabla_{\partial K}u,\nabla_{\partial K}u)\,
      \dd\tilde{\mu}_{\alpha,\partial K}
    -
    2\int_{\partial K}
      \big\langle
        \nabla_{\partial K}\varphi,
        \nabla_{\partial K}u
      \big\rangle
      \dd\tilde{\mu}_{\alpha,\partial K}  \\
    &\ge
    \int_{\partial K}
      H_{W_\alpha}\varphi^2\,
      \dd\tilde{\mu}_{\alpha,\partial K}
    +
    \int_{\partial K}
      \mathrm{II}(\nabla_{\partial K}u,\nabla_{\partial K}u)\,
      \dd\tilde{\mu}_{\alpha,\partial K}       \\
    &\quad
    -
    2\int_{\partial K}
      \big\langle
        \nabla_{\partial K}\varphi,
        \nabla_{\partial K}u
      \big\rangle
      \dd\tilde{\mu}_{\alpha,\partial K}.
  \end{align*}
  Using $\mathrm{II}>0$ for $K\in\mathcal{K}_{+}^n$, and applying the Cauchy-Schwarz inequality, we deduce
    \[
    \frac{\tilde{\mu}'(0)^2}{\tilde{\mu}(0)}
    \ge
    \int_{\partial K}
    \left(
      H_{W_\alpha}\varphi^2
      -
      \mathrm{II}^{-1}
      \bigl(
        \nabla_{\partial K}\varphi,
        \nabla_{\partial K}\varphi
      \bigr)
    \right)
    \dd\tilde{\mu}_{\alpha,\partial K}
    =\tilde{\mu}''(0).
  \]
  Moreover, if $\psi\not\equiv0$ and the inequality is not strict, then the equality implies $I_\alpha(u)=0$. Proposition~\ref{prop-energy} then forces $\varphi = u_\nu = 0$ on $\partial K$. Since $\varphi = \psi \circ \nu_K$ and the Gauss map $\nu_K$ is onto $\mathbb{S}^{n-1}$, we obtain $\psi \equiv 0$, a contradiction. Hence the inequality is strict.
\end{proof}

\end{document}
